\label{sec:related}
\textbf{AI coding agents on GitHub.} Research on coding agents has moved quickly from benchmarks to field data. SWE-bench measures whether a model can resolve real GitHub issues~\cite{jimenez2024swebench}, and SWE-agent showed that a purpose-built interface lets a language model navigate a repository, edit files and run tests to do so~\cite{yang2024sweagent}. Field studies now follow agents into real projects. Li et al.\ describe the arrival of autonomous agents as ``AI teammates'' and released AIDev, a large dataset of agent-authored pull requests; agent pull requests are accepted less often than human ones~\cite{li2025aiteammates,li2026aidev}. Pinna et al.\ compared five agents and found that acceptance depends strongly on the type of task, so that a single global ranking of agents can mislead~\cite{pinna2026comparing}. Watanabe et al.\ studied pull requests made with Claude Code and found that developers often delegate refactoring, documentation and testing to it~\cite{watanabe2025agentic}. Pham and Ghaleb compared the code changes and descriptions of agentic and human pull requests~\cite{pham2026alignment}, and Selvanayagam and Ghaleb studied pull requests in which one agent reviews another agent's code~\cite{selvanayagam2026aitoai}. These studies focus on pull requests. Sayagh is an exception: the study examines which GitHub issues lead to merged pull requests by the Copilot agent~\cite{sayagh2025copilot}, but it treats the issue as input for one agent. The present study looks instead at the issues attributed to thirteen agents, whatever role the agent plays in them.

\textbf{Attributing activity to agents and bots.} Knowing who or what produced an artefact is a research problem of its own. Bot accounts are common on GitHub: Wessel et al.\ found bots in 26\% of the popular projects they sampled~\cite{wessel2018bots}, and Golzadeh et al.\ built a ground-truth dataset and classifier to separate bot from human accounts in issue and pull-request comments~\cite{golzadeh2021bots}. For coding agents the problem is harder, because agent output is often posted under human accounts. Ghaleb showed that agents leave fingerprints that a classifier can use to identify them~\cite{ghaleb2026fingerprinting}, and the CodAGE dataset, from which the data of this study appear to derive, assigns agents with login, branch and text signals and calls the result a candidate attribution with false positives by design~\cite{ghaleb2026codage}. Kalliamvakou et al.\ warned more generally that GitHub data contain many such perils for researchers~\cite{kalliamvakou2014promises}. The analysis here carries that warning down to single issues: it records how each issue is linked to its agent, and the link turns out to differ sharply between agents.

\textbf{Bots in maintainers' work.} Bots change how projects work. When projects adopted code-review bots, more pull requests were merged and developers communicated less~\cite{wessel2020effects}. Practitioners nevertheless disagree about what counts as a development bot and what it is good for~\cite{erlenhov2020bots}, and developers describe unwanted bot activity as noise that disrupts their work~\cite{wessel2021disturb}. Agents that open their own issues add to this activity. Whether maintainers act on those issues has, as far as the researcher is aware, not been measured; RQ1 and RQ2 do so.

\textbf{Issue labels, lifetime and topics.} Labels are GitHub's built-in way to classify and triage issues. Their use varies between projects and is associated with how issues are resolved~\cite{cabot2015labels}, and tools such as Ticket Tagger assign them automatically~\cite{kallis2019tagger}. Labels and issue types are noisy, however: Herzig et al.\ inspected more than 7{,}000 issue reports by hand and found that about a third were misclassified~\cite{herzig2013bug}. The lifetime of an issue has been studied as a prediction target~\cite{kikas2016lifetime}. To learn what issues are about, mining studies most often use LDA~\cite{blei2003lda,chen2016survey,silva2021topic}, although its topics can change with the random seed~\cite{agrawal2018wrong,greene2014stability}. Recent studies of agent pull requests and of ChatGPT use in issue trackers use BERTopic instead~\cite{opu2025performance,das2024chatgpt,grootendorst2022bertopic}. Labels, lifetime and topics are combined here; open issues are handled with survival analysis, and topics are reported only after their stability has been tested and a second topic model has confirmed them.
